\section{Introduction}

We set out to measure whether better pre-training data produces more generalizable language models --- and found the opposite. OLMo-7B, trained on Dolma~\cite{Soldaini2024Dolma}, one of the most carefully curated pre-training corpora to date with multi-stage quality filtering and explicit inclusion of academic content, achieves a lower knowledge-to-commonsense generalization balance ratio than Pythia-6.9B~\cite{Biderman2023Pythia}, trained on The Pile with minimal curation. More strikingly, both models converge to \textit{identical} commonsense reasoning scores, suggesting that web-sourced commonsense knowledge saturates at similar levels regardless of curation quality at this training scale.

The intuition that better-curated data produces better-generalized models drives enormous resource investment across industry and academia. Dolma's curators deduplicated, filtered for quality, and explicitly included academic sources (S2ORC, Wikipedia) precisely because these choices were expected to yield models with stronger knowledge acquisition and more robust generalization. Yet at approximately 300 billion training tokens --- a scale at which both OLMo-7B and Pythia-6.9B have been released with documented checkpoints --- the model trained on the less-curated corpus achieves a higher MMLU/HellaSwag generalization balance ratio (Pythia: 0.565 vs OLMo: 0.538; difference $= -0.0265$, 95\% CI $[-0.045, -0.007]$, Cohen's $d = -2.732$).

\textbf{The deeper problem.} This result cannot be cleanly attributed to corpus quality alone. OLMo-7B and Pythia-6.9B differ in architecture, tokenizer, and optimizer configuration~\cite{Groeneveld2024OLMo,Biderman2023Pythia}. Any performance difference is therefore a joint function of these factors, and disentangling corpus quality from architecture requires either matched-architecture experiments or training trajectory analysis across multiple checkpoints. Our study is not that experiment --- but it reveals, precisely, why that experiment is necessary.

\textbf{The gap.} Prior work on corpus curation~\cite{Penedo2023RefinedWeb,Soldaini2024Dolma,Muennighoff2023Scaling} has demonstrated the importance of data quality in aggregate, but controlled comparisons at matched training scale and matched architecture are rare. Most corpus comparisons are confounded by training duration, architecture differences, or evaluation inconsistencies. The claim that curation quality drives generalization has largely gone untested under the conditions required to isolate it.

\textbf{Our finding.} At $\sim$300B tokens, the curated-corpus hypothesis is not supported: Pythia-6.9B achieves a higher MMLU/HellaSwag ratio than OLMo-7B. The most structurally informative aspect is that HellaSwag scores are identical (0.4580 each), suggesting denominator saturation at this scale --- which means the ratio is driven entirely by MMLU differences that are themselves architecture-confounded. This null (and negative) result has direct implications for how corpus quality studies should be designed and how generalization balance metrics should be validated before use.

\subsection*{Contributions}

\begin{enumerate}
    \item \textbf{Controlled matched-scale evaluation.} We provide the first systematic comparison of OLMo-7B and Pythia-6.9B at precisely matched training tokens ($\sim$300B), with documented checkpoint provenance and reproducible evaluation via \texttt{lm-evaluation-harness}.
    \item \textbf{Refutation of the curated-corpus hypothesis at this scale.} We demonstrate that OLMo-7B does not achieve a higher MMLU/HellaSwag generalization balance ratio than Pythia-6.9B, with a statistically robust reversal (95\% CI entirely below zero, Cohen's $d = -2.732$).
    \item \textbf{Structural finding: HellaSwag denominator saturation.} We identify that identical HellaSwag scores (0.4580 each) indicate saturation of the commonsense benchmark at 6--8B model scale and $\sim$300B tokens, collapsing ratio metrics into MMLU proxies and limiting the interpretability of generalization balance comparisons.
    \item \textbf{Methodological recommendations.} We characterize the minimum design requirements for future corpus quality studies: architecture-matched baselines or trajectory analysis, denominator-sensitivity checks before ratio metric adoption, and contamination controls for MMLU-adjacent benchmarks.
\end{enumerate}

\subsection*{Paper Organization}

Section~2 reviews related work on corpus curation, generalization metrics, and data attribution. Section~3 describes our methodology, including checkpoint selection, evaluation protocol, and statistical analysis. Section~4 details the experimental setup. Section~5 presents results. Section~6 discusses interpretations, limitations, and implications for metric design. Section~7 concludes with future directions.
